\chapter{Decision trees and recursive partitions}
\label{ch:decision-trees}

\begin{learningobjectives}
After this chapter, you should be able to:
\begin{itemize}
  \item interpret a decision tree as a recursive partition and prediction rule;
  \item compare regression and classification split criteria;
  \item control complexity through structural constraints and pruning; and
  \item audit threshold, stability, missingness, and interpretation failures.
\end{itemize}
\end{learningobjectives}

\begin{chapterconnection}
Geometric methods defined neighborhoods and margins globally. A decision tree
constructs axis-aligned local regions through a sequence of binary questions.
This representation naturally captures nonlinearities and interactions, but its
greedy construction is unstable. That instability motivates the ensembles in
the next chapter.
\end{chapterconnection}

\section{A tree is a piecewise-constant function}

An internal node asks whether $x_j\le t$. The answer routes an observation to a
child. A leaf returns a constant: commonly a target mean for regression or class
frequencies for classification. The root-to-leaf path defines a rectangular
region of feature space.

Tree fitting greedily chooses a feature and threshold that reduce an impurity
criterion. For regression, reduction in within-node squared error is common. For
classification with proportions $p_k$, Gini impurity is
$1-\sum_kp_k^2$ and entropy is $-\sum_kp_k\log p_k$. These criteria guide local
splits; they are not the final deployment utility.

\section{Greedy growth and structural regularization}

Searching every possible tree is combinatorial. Standard algorithms choose the
best immediate split, partition the data, and repeat. A locally best split need
not belong to the globally best tree.

Depth, minimum samples per split or leaf, maximum leaf count, and minimum
impurity decrease constrain growth. Cost-complexity pruning instead grows a
larger tree and trades empirical fit against leaf count. These are forms of
regularization and must be chosen within the validation boundary.

\conceptcheck{Why can a leaf with one training observation have zero regression
impurity and nevertheless be a poor estimate of its region's conditional mean?}

\section{Instability is a central property}

A small data perturbation can change an early split and therefore the entire
downstream partition. Two trees can have similar predictive performance but
very different structures. A displayed tree is an exact account of one fitted
artifact, not necessarily a stable description of the population.

Measure stability by refitting across resamples or folds and comparing
predictions, split variables, and subgroup behavior. The next chapter will use
this variance as an opportunity: averaging decorrelated trees can produce a
much more stable predictor.

\section{Threshold semantics and missing values}

Tree boundaries create discontinuities. Inputs just below and above a threshold
may receive sharply different predictions. Tests should exercise equality and
neighboring representable values. Unit conversions and numerical rounding can
therefore change paths.

Missingness needs a declared policy: rejection, imputation inside the pipeline,
learned default direction, or a native missing-value mechanism. Missingness may
itself encode collection behavior that shifts after deployment. Category
handling also needs stable vocabularies and an unknown-category policy.

\begin{professionalbox}
Export the exact feature schema, preprocessing, library and artifact versions,
class ordering, and missing-value behavior with a tree. Test a small set of
golden predictions at every release, including values at important split
thresholds.
\end{professionalbox}

\section{Interpretability has a scope}

A path answers: ``Which conditions in this fitted tree produced this output?''
It does not answer why the outcome occurred, what an intervention would do, or
whether the same path would appear under resampling. Impurity-based feature
importance can favor variables with many candidate splits and distribute credit
unreliably among correlated variables.

Permutation and held-out perturbation methods ask different questions and can
also fail under dependence or unrealistic perturbations. Interpretation should
name the model, population, perturbation, and uncertainty.

\begin{misconceptionbox}
\textbf{Misconception:} a shallow tree is inherently objective and
interpretable. \textbf{Correction:} it may be readable, but its variables,
thresholds, sample, missingness policy, and pruning objective still encode
choices. Readability does not establish stability, fairness, or causality.
\end{misconceptionbox}

\begin{workedexample}{a brittle degradation rule}
A depth-three tree places resistance change at its root and performs well on a
held-out random split. When the evaluation holds out later manufacturing
batches, performance falls. Repeated fits move the root threshold substantially.

Investigation shows that the instrument calibration changed between batches.
The tree had learned a convenient separator tied to measurement context. The
team retains the tree as a baseline, documents its threshold instability, adds
batch-aware validation, and rejects causal language about resistance. The
failure is useful evidence about the data-generating system.
\end{workedexample}

\section{Exercises}
\begin{enumerate}
  \item Compute Gini impurity and entropy before and after a proposed binary
  split. Compare their ordering.
  \item Fit the conceptual structure of two trees to bootstrap samples and list
  the quantities needed for a stability report.
  \item Design boundary tests for a threshold at $0.2$ with missing, infinite,
  and unit-converted inputs.
  \item Explain why impurity importance is not a causal effect and not an
  automatic measure of predictive necessity.
\end{enumerate}

\begin{chapterreview}
Decision trees construct piecewise-constant predictions through greedy
recursive partitions. Split criteria, structural constraints, and pruning
define the learner. Trees represent nonlinear interactions naturally, but they
are discontinuous and often unstable. Their readable paths explain one artifact
only within a carefully stated scope.
\end{chapterreview}

\begin{notebooklink}
Lecture 13 notebook 00 develops recursive partitioning and impurity. Its second
laboratory studies pruning, sample perturbation, threshold tests, and the limits
of tree interpretation.
\end{notebooklink}
